\begin{thebibliography}{9}
\addcontentsline{toc}{chapter}{References}
\markboth{References}{}
\bibitem{Montalban}
Antonio Montalb\'an.
\emph{Computable Structure Theory: Within the Arithmetic}.
Perspectives in Logic, Cambridge University Press, 2021.
For the cited numbering, see the author's draft of June 25, 2019, Example III.39 and the following paragraph, p.~50.
\url{https://math.berkeley.edu/~antonio/CSTpart1.pdf}.

\bibitem{Calvert}
Wesley Calvert, Valentina Harizanov, and Alexandra Shlapentokh.
Turing degrees of isomorphism types of algebraic objects.
\emph{Journal of the London Mathematical Society} 75(2) (2007), 273--286.
\href{https://doi.org/10.1112/jlms/jdl012}{doi:10.1112/jlms/jdl012}.
Preprint: \href{https://arxiv.org/abs/math/0507128}{arXiv:math/0507128}.

\bibitem{Fuchs}
L\'aszl\'o Fuchs.
\emph{Abelian Groups}.
Springer Monographs in Mathematics, Springer, 2015.
\href{https://doi.org/10.1007/978-3-319-19422-6}{doi:10.1007/978-3-319-19422-6}.
General background for the terminology of pure subgroups and mixed abelian groups. The arguments used here are proved in the text.

\bibitem{Lupini}
Martino Lupini, Alexander Melnikov, and Andr\'e Nies.
\emph{Computable topological abelian groups}.
\href{https://arxiv.org/abs/2105.12897v2}{arXiv:2105.12897v2}, 2021.
Definition 3.3 supplies the computable inverse-system convention. Section 3.2 describes symmetric cocycles and abelian extensions.

\bibitem{Wilkes}
Gareth Wilkes.
\emph{Part III Profinite Groups}.
Lecture notes, Lent Term 2020.
Sections 1.2 and 3.1 discuss profinite completion and its functoriality.
\url{https://www.dpmms.cam.ac.uk/~grw46/LectureNotes.pdf}.
\end{thebibliography}
